\chapter{वर्गीकरण और नातेदारी}\label{ch:g6:classification-kinship}

\cref{ch:g4:classifying-animals} का \omterm{def:g4:classifying-animals:classification}{वर्गीकरण}
इतना अच्छा काम क्यों करता है? \omterm{def:g1:animals-around-us:animal}{जंतुओं} को \omterm{def:g1:animals-around-us:covering}{पंखों} और \omterm{def:g1:animals-around-us:covering}{रोमों} से छाँटो, और
हज़ार दूसरी ख़ासियतें आज्ञाकारी होकर क़तार में आ जाती हैं --- \omterm{def:g2:animals-grow-up:born}{दूध}
पिलाने वालों के पास \omterm{def:g4:classifying-animals:vertebrate}{रीढ़} होती है, \omterm{def:g1:animals-around-us:covering}{पंख} पहनने वाले खोलदार \omterm{def:g2:animals-grow-up:born}{अंडे} देते हैं,
और किसी को जाँचना तक नहीं पड़ता।
\Cref{rem:g4:classifying-animals:tree} ने यह सवाल अधूरा छोड़ दिया था।
यह अध्याय उसका उत्तर एक ही शब्द से देता है --- \omterm{prop:g6:classification-kinship:kinship}{नातेदारी} --- और उन
डिब्बों को कुल-वृक्ष के रूप में फिर से खड़ा करता है।

\section{अभिलक्षण, साझे और नत्थी}

\begin{definition}[अभिलक्षण]\label{def:g6:classification-kinship:attribute}
\emph{अभिलक्षण}\index{अभिलक्षण} किसी जीव की वह ख़ासियत है जो उसके
\emph{पास होती है}, और जो इस तरह कही जाती है कि जाँची जा सके: \omterm{def:g4:classifying-animals:vertebrate}{रीढ़} वाला
\omterm{def:g3:bones-and-muscles:skeleton}{कंकाल} है; चार \omterm{def:g1:my-body:parts}{उपांग} हैं; \omterm{def:g1:animals-around-us:covering}{रोम} और \omterm{def:g2:animals-grow-up:born}{दूध} पिलाना है; \omterm{def:g1:animals-around-us:covering}{पंख} हैं। \omterm{def:g4:classifying-animals:classification}{वर्गीकरण}
अभिलक्षणों पर ही खड़ा है --- यानी
\cref{def:g3:sorting-living-things:criterion} के अच्छे \omterm{def:g3:sorting-living-things:criterion}{मानदंड}, क़ायदे
में ढले हुए।
\end{definition}

\begin{proposition}[अभिलक्षण नत्थी होते हैं]\label{prop:g6:classification-kinship:nest}
बहुत-सी \omterm{def:g5:biodiversity:species}{प्रजातियों} पर जाँचने पर \omterm{def:g6:classification-kinship:attribute}{अभिलक्षण} \emph{नत्थी} समूहों में बैठ
जाते हैं --- \omterm{def:g1:animals-around-us:covering}{रोम} वाले हर \omterm{def:g1:animals-around-us:animal}{जंतु} के चार \omterm{def:g1:my-body:parts}{उपांग} भी होते हैं; चार उपांगों
वाले हर \omterm{def:g1:animals-around-us:animal}{जंतु} के \omterm{def:g4:classifying-animals:vertebrate}{रीढ़} भी होती है; और उलटा ढर्रा कभी नहीं मिलता (\omterm{def:g4:classifying-animals:vertebrate}{रीढ़} के
बिना \omterm{def:g1:animals-around-us:covering}{रोम}) और न समूह कभी आपस में कटते हैं (\omterm{def:g1:animals-around-us:covering}{रोम} और \omterm{def:g1:animals-around-us:covering}{पंख} एक साथ)। यही
नत्थी होना वह वजह है कि समूहों के भीतर समूह
(\cref{def:g4:classifying-animals:classification}) जीवित दुनिया पर
बैठते ही क्यों हैं।
\end{proposition}
\admitted

\begin{omfigure}
\begin{tikzpicture}[scale=0.9]
  % nested attribute boxes
  \draw[thick, omDef, fill=omDef!4] (0,0) rectangle (11.4,4.4);
  \node[font=\small\bfseries, omDef] at (2.2,4.1)
    {रीढ़ है};
  \draw[thick, omMeth, fill=omMeth!5] (0.4,0.35) rectangle (8.6,3.5);
  \node[font=\small\bfseries, omMeth] at (2.35,3.2)
    {चार उपांग हैं};
  \draw[thick, omProp, fill=omProp!6] (0.8,0.7) rectangle (5.6,2.6);
  \node[font=\small\bfseries, omProp] at (2.6,2.3)
    {रोम और दूध हैं};
  \node[font=\small] at (3.2,1.4) {बिल्ली, चमगादड़, ह्वेल};
  \node[font=\small, align=center] at (7.0,1.6)
    {मेंढक, छिपकली,\\ कलचिड़ी};
  \node[font=\small] at (10.0,1.6) {ट्राउट};
\end{tikzpicture}

{\small \omterm{def:g6:classification-kinship:attribute}{अभिलक्षण} नत्थी होते हैं: \omterm{def:g1:animals-around-us:covering}{रोम} वाले चार उपांगों वालों के भीतर
बैठते हैं, और चार उपांगों वाले \omterm{def:g4:classifying-animals:vertebrate}{रीढ़} वालों के भीतर। ट्राउट के पास \omterm{def:g4:classifying-animals:vertebrate}{रीढ़}
तो है पर वह चार उपांगों वाले डिब्बे से बाहर खड़ी है।}
\end{omfigure}

\begin{example}[तीन जंतुओं को रखना]\label{ex:g6:classification-kinship:placing}
ह्वेल, ट्राउट और कलचिड़ी के लिए डिब्बे जाँचो। ह्वेल: \omterm{def:g4:classifying-animals:vertebrate}{रीढ़} हाँ, चार \omterm{def:g1:my-body:parts}{उपांग}
हाँ (आगे मीन-पंख, और पीछे छिपे हुए अवशेष), रोम-और-दूध हाँ --- यानी
सबसे भीतरी डिब्बा, बिल्ली के साथ। ट्राउट: \omterm{def:g4:classifying-animals:vertebrate}{रीढ़} हाँ, बस --- सिर्फ़ बाहरी
डिब्बा। कलचिड़ी: \omterm{def:g4:classifying-animals:vertebrate}{रीढ़}, चार \omterm{def:g1:my-body:parts}{उपांग} (दो \omterm{def:g1:my-body:parts}{टाँगें}, दो डैने), \omterm{def:g1:animals-around-us:covering}{पंख} --- यानी चार
उपांगों वाला डिब्बा, और उसमें \omterm{def:g1:animals-around-us:covering}{पंखों} का अपना ख़ाना। कोई \omterm{def:g1:animals-around-us:animal}{जंतु} कभी एक साथ
दो कटते हुए डिब्बे नहीं माँगता।
\end{example}

\section{नातेदारी: नत्थी होने के पीछे का कारण}

\begin{proposition}[साझे अभिलक्षण साझे पूर्वजों का मतलब हैं]\label{prop:g6:classification-kinship:kinship}
इस नत्थी होने की व्याख्या वंशागति है
(\cref{rem:g4:animal-reproduction:mix}: बच्चे अपने जनकों से पाते हैं)।
कोई \omterm{def:g6:classification-kinship:attribute}{अभिलक्षण} साझा करने वाली \omterm{def:g5:biodiversity:species}{प्रजातियाँ} उसे इसलिए साझा करती हैं कि
उन्होंने उसे \emph{किसी साझे पूर्वज से विरासत में पाया} --- यानी किसी
दूर की दादी-प्रजाति से। दो \omterm{def:g5:biodiversity:species}{प्रजातियाँ} जितने ज़्यादा \omterm{def:g6:classification-kinship:attribute}{अभिलक्षण} साझा करती
हैं, उनकी \emph{नातेदारी}\index{नातेदारी} उतनी ही क़रीबी है: बिल्ली और
ह्वेल, दोनों \omterm{def:g1:animals-around-us:covering}{रोम} वाली \omterm{def:g2:animals-grow-up:born}{दूध} पिलाने वाली, एक-दूसरे की उससे ज़्यादा क़रीबी
चचेरी हैं जितनी उनमें से कोई ट्राउट की है। \omterm{def:g4:classifying-animals:classification}{वर्गीकरण} का हर समूह असल में
एक \emph{कुल} है: कोई पूर्वज और उसकी सारी संतानें।
\end{proposition}
\admitted

\begin{example}[ह्वेल, आख़िरकार ठीक से पढ़ी हुई]\label{ex:g6:classification-kinship:whale}
इसीलिए ह्वेल वाली पहेली
(\cref{exo:g4:classifying-animals:9}) को वैसे ही सुलझना था जैसे वह
सुलझी: रोम-और-दूध ह्वेल की कुल-रेखा पर निशान लगाते हैं --- उसके पूर्वज
\omterm{def:g2:animals-grow-up:born}{दूध} पिलाने वाले थलचर \omterm{prop:g3:sorting-living-things:groups}{स्तनधारी} थे --- जबकि उसकी मछली जैसी बनावट पानी की
करनी है (\cref{exo:g4:adapted-to-their-world:11}), जो एक असंबंधित कुल
पर उसी जीवन-ढंग ने छाप दी। पूर्वजों वाले \omterm{def:g6:classification-kinship:attribute}{अभिलक्षण} कुलों को छाँटते हैं;
और \omterm{def:g2:where-animals-live:habitat}{आवास} की समानताएँ अजनबी भी उधार ले सकते हैं। \omterm{def:g4:classifying-animals:classification}{वर्गीकरण} पहले पर भरोसा
करता है और दूसरे पर शक।
\end{example}

\begin{definition}[नातेदारी-वृक्ष]\label{def:g6:classification-kinship:tree}
\omterm{prop:g6:classification-kinship:kinship}{नातेदारी} को एक \emph{वृक्ष}\index{नातेदारी-वृक्ष} की तरह खींचा जाता है:
शाखाओं के बिंदु साझे पूर्वजों को दिखाते हैं, और हर \omterm{def:g6:classification-kinship:attribute}{अभिलक्षण} एक ही बार
आता है, उसी शाखा पर जहाँ वह पैदा हुआ था --- और उस दोराहे से आगे की हर
\omterm{def:g5:biodiversity:species}{प्रजाति} उसे विरासत में पाती है। नत्थी डिब्बे और यह वृक्ष एक ही जानकारी
लिए रहते हैं: डिब्बा वही शाखा है, अपनी सारी टहनियों समेत।
\end{definition}

\begin{omfigure}
\begin{tikzpicture}[scale=0.95]
  % tree
  \draw[very thick] (0,0) -- (2.2,0);
  % backbone mark
  \fill[omDef] (1.6,0) circle (0.09);
  \node[font=\footnotesize, omDef, align=center] at (1.6,-0.55)
    {रीढ़ का\\ प्रकट होना};
  \draw[very thick] (2.2,0) -- (3.6,1.6) -- (10.4,1.6);
  \draw[very thick] (2.2,0) -- (4.4,-1.4) -- (10.4,-1.4);
  \node[font=\small, right] at (10.5,-1.4) {ट्राउट};
  % four limbs mark
  \fill[omMeth] (3.1,1.05) circle (0.09);
  \node[font=\footnotesize, omMeth, align=center] at (2.4,1.45)
    {चार उपांगों का\\ प्रकट होना};
  \draw[very thick] (5.0,1.6) -- (6.2,2.8) -- (10.4,2.8);
  % fur mark on upper branch
  \fill[omProp] (5.75,2.35) circle (0.09);
  \node[font=\footnotesize, omProp, align=center] at (5.0,2.9)
    {रोम और दूध का\\ प्रकट होना};
  \draw[very thick] (6.9,2.8) -- (7.7,3.7) -- (10.4,3.7);
  \node[font=\small, right] at (10.5,3.7) {बिल्ली};
  \node[font=\small, right] at (10.5,2.8) {ह्वेल};
  % feathers on middle branch
  \fill[omThm] (7.4,1.6) circle (0.09);
  \node[font=\footnotesize, omThm, align=center] at (7.4,1.05)
    {पंखों का\\ प्रकट होना};
  \node[font=\small, right] at (10.5,1.6) {कलचिड़ी};
\end{tikzpicture}

{\small एक जेबी \omterm{def:g6:classification-kinship:tree}{नातेदारी-वृक्ष}। हर रंगीन बिंदु एक \omterm{def:g6:classification-kinship:attribute}{अभिलक्षण} है, जो एक ही
बार आता है; और उससे आगे की हर शाखा उसे विरासत में पाती है। बिल्ली और
ह्वेल सबसे नया दोराहा साझा करती हैं --- यानी यहाँ खींचे गए सबसे क़रीबी
चचेरे।}
\end{omfigure}

\begin{method}[छोटा नातेदारी-वृक्ष बनाना]\label{met:g6:classification-kinship:build}
मुट्ठी भर \omterm{def:g5:biodiversity:species}{प्रजातियाँ} सामने हों, तो:
\begin{enumerate}
  \item उनके \omterm{def:g6:classification-kinship:attribute}{अभिलक्षणों} की तालिका बनाओ --- पंक्तियों में \omterm{def:g5:biodiversity:species}{प्रजातियाँ},
        स्तंभों में जाँचे जा सकने वाले \omterm{def:g6:classification-kinship:attribute}{अभिलक्षण}, और जहाँ मौजूद हों वहाँ
        निशान;
  \item साझे निशानों से समूह बनाओ, सबसे चौड़े \omterm{def:g6:classification-kinship:attribute}{अभिलक्षण} से शुरू करके:
        तालिका में ही नत्थी होना दिख जाता है;
  \item \omterm{def:g6:classification-kinship:tree}{वृक्ष} खींचो: हर \omterm{def:g6:classification-kinship:attribute}{अभिलक्षण} के लिए एक दोराहा, और सबसे चौड़ा तने के
        सबसे पास; हर \omterm{def:g5:biodiversity:species}{प्रजाति} को ठीक उन्हीं दोराहों से आगे रखो जिनके
        \omterm{def:g6:classification-kinship:attribute}{अभिलक्षण} उसके पास हैं;
  \item \omterm{prop:g6:classification-kinship:kinship}{नातेदारी} पढ़ो: दो \omterm{def:g5:biodiversity:species}{प्रजातियों} की क़रीबी वही आख़िरी दोराहा है जो
        वे साझा करती हैं।
\end{enumerate}
\end{method}

\begin{remark}[वृक्ष चुपचाप क्या कह देता है]\label{rem:g6:classification-kinship:implies}
इस \omterm{def:g6:classification-kinship:tree}{वृक्ष} को गंभीरता से लो और वह कुछ बहुत बड़ा कह देता है: किन्हीं भी दो
शाखाओं का पीछे तक पीछा करो और वे \emph{मिल जाती} हैं। बिल्ली और ट्राउट
की एक दादी-प्रजाति साझी है; तो और पीछे जाकर बिल्ली और घोंघे की भी होनी
चाहिए, और गुलबहार तथा बच्चे की भी। तब तो सारे \omterm{def:g6:exploring-our-environment:components}{सजीव} एक ही कुल हुए, और
उनकी एकता (\cref{prop:g5:first-look-at-cells:principle}) एक कुल की
पारिवारिक समानता --- और तब \omterm{def:g5:biodiversity:species}{प्रजातियों} का पीढ़ियों में \emph{बदल} सकना
भी ज़रूरी है, वरना यह शाखाएँ बनतीं ही नहीं। यह सच है या नहीं, और यह हो
कैसे सकता है, यही \cref{ch:g9:evidence-of-evolution} वाले साल का बड़ा
सवाल है --- और यह \omterm{def:g6:classification-kinship:tree}{वृक्ष} वही चित्र है जो उसे पूछने पर मजबूर कर देता है।
\end{remark}

\section{अभ्यास}

\begin{exercise}[$\star$]\label{exo:g6:classification-kinship:1}
\omterm{def:g6:classification-kinship:attribute}{अभिलक्षण} क्या है? तीन बताओ, और जाँचे जा सकने लायक़ ढंग से।
\end{exercise}

\begin{exercise}[$\star$]\label{exo:g6:classification-kinship:2}
\omterm{def:g6:classification-kinship:attribute}{अभिलक्षणों} का \emph{नत्थी} होना किसे कहते हैं? अध्याय का तीन डिब्बों
वाला उदाहरण दो।
\end{exercise}

\begin{exercise}[$\star$]\label{exo:g6:classification-kinship:3}
\cref{prop:g6:classification-kinship:kinship} के अनुसार \omterm{def:g5:biodiversity:species}{प्रजातियाँ}
\omterm{def:g6:classification-kinship:attribute}{अभिलक्षण} साझा क्यों करती हैं?
\end{exercise}

\begin{exercise}[$\star$]\label{exo:g6:classification-kinship:4}
जेबी \omterm{def:g6:classification-kinship:tree}{वृक्ष} पर ट्राउट किस अभिलक्षण-बिंदु से आगे खड़ी है? और ह्वेल किन
बिंदुओं से आगे?
\end{exercise}

\begin{exercise}[$\star$]\label{exo:g6:classification-kinship:5}
ज़्यादा क़रीबी चचेरे कौन हैं: ह्वेल और बिल्ली, या ह्वेल और ट्राउट?
उत्तर \omterm{def:g6:classification-kinship:tree}{वृक्ष} के दोराहों से पढ़ो।
\end{exercise}

\begin{exercise}[$\star$]\label{exo:g6:classification-kinship:6}
किसी नत्थी डिब्बे और \omterm{def:g6:classification-kinship:tree}{वृक्ष} की किसी शाखा का संबंध बताओ।
\end{exercise}

\begin{exercise}[$\star\star$]\label{exo:g6:classification-kinship:7}
\omterm{def:g4:classifying-animals:classification}{वर्गीकरण} रोम-और-दूध पर भरोसा और ``मछली जैसी बनावट'' पर शक क्यों करता
है? \cref{ex:g6:classification-kinship:whale} का इस्तेमाल करो।
\end{exercise}

\begin{exercise}[$\star\star$]\label{exo:g6:classification-kinship:8}
\cref{met:g6:classification-kinship:build} की अभिलक्षण-तालिका इनके लिए
बनाओ: ट्राउट, मेंढक, बिल्ली, कलचिड़ी --- और \omterm{def:g6:classification-kinship:attribute}{अभिलक्षण} ये हों: \omterm{def:g4:classifying-animals:vertebrate}{रीढ़}, चार
\omterm{def:g1:my-body:parts}{उपांग}, रोम-और-दूध, \omterm{def:g1:animals-around-us:covering}{पंख}।
\end{exercise}

\begin{exercise}[$\star\star$]\label{exo:g6:classification-kinship:9}
जेबी \omterm{def:g6:classification-kinship:tree}{वृक्ष} में छिपकली जोड़ो: वह किन दोराहों से गुज़रती है, और उसकी शाखा
कहाँ अलग होती है?
\end{exercise}

\begin{exercise}[$\star\star$]\label{exo:g6:classification-kinship:10}
चमगादड़ कलचिड़ी की तरह उड़ता है पर बिल्ली की तरह \omterm{def:g2:animals-grow-up:born}{दूध} पिलाता है। कौन-सी
समानता \omterm{prop:g6:classification-kinship:kinship}{नातेदारी} बताती है और कौन-सी जीवन-ढंग से उधार ली गई है? \omterm{def:g6:classification-kinship:tree}{वृक्ष} के
तर्क से तय करो।
\end{exercise}

\begin{exercise}[$\star\star$]\label{exo:g6:classification-kinship:11}
``कोई समूह एक पूर्वज और उसकी सारी संतानें होता है।'' इस वाक्य को
\omterm{prop:g3:sorting-living-things:groups}{स्तनधारियों} पर और \cref{exo:g4:classifying-animals:11} के प्लैटिपस पर
आज़माओ: क्या \omterm{def:g2:animals-grow-up:born}{अंडे} देने वाली वह अनोखी बात कुल तोड़ देती है, या उसी की
है?
\end{exercise}

\begin{exercise}[$\star\star\star$]\label{exo:g6:classification-kinship:12}
\cref{rem:g6:classification-kinship:implies} को उसके सिरे तक ले जाओ:
अगर शाखाओं की हर जोड़ी सचमुच मिलती है, तो लंबे समय में \omterm{def:g5:biodiversity:species}{प्रजातियों} के
बारे में क्या सच होना चाहिए? इस निहितार्थ को एक वाक्य में लिखो --- और
यह भी बताओ कि तुम इसके लिए किस तरह का सबूत माँगोगे।
\end{exercise}

\section{समस्या: संग्रहालय की दराज़}

\begin{problem}[{सप्ताहांत समस्या --- छह नमूने, एक अभिलक्षण-तालिका, और
एक वृक्ष}]\label{pb:g6:classification-kinship:1}
एक संग्रहालय कक्षा को छह नमूनों की दराज़ उधार देता है: एक ट्राउट, एक
मेंढक, एक छिपकली, एक कबूतर, एक चूहा, और मीन-पंखों के ढाँचों समेत
डॉल्फ़िन की एक \omterm{def:g3:bones-and-muscles:skeleton}{खोपड़ी}। काम यह है: इन्हें \omterm{def:g6:classification-kinship:attribute}{अभिलक्षणों} से वर्गीकृत करो और
इनका \omterm{def:g6:classification-kinship:tree}{नातेदारी-वृक्ष} खींचो।

\medskip\noindent\textbf{भाग I --- तालिका।}
\begin{enumerate}
  \item छहों के लिए तालिका बनाओ, और \omterm{def:g6:classification-kinship:attribute}{अभिलक्षण} ये लो: \omterm{def:g4:classifying-animals:vertebrate}{रीढ़}; चार \omterm{def:g1:my-body:parts}{उपांग};
        सूखी शल्कदार \omterm{def:g1:the-five-senses:sense}{त्वचा}; \omterm{def:g1:animals-around-us:covering}{पंख}; और \omterm{def:g1:animals-around-us:covering}{रोम} तथा \omterm{def:g2:animals-grow-up:born}{दूध}। (डॉल्फ़िन के लगभग कोई
        बाल नहीं होते, पर संग्रहालय की चिट गवाही देती है: बच्चे \omterm{def:g2:animals-grow-up:born}{दूध} पर
        पलते हैं --- \cref{ex:g3:sorting-living-things:surprises}
        समझाता है कि यह गवाही मामला क्यों तय कर देती है।)
  \item छहों में कौन-सा \omterm{def:g6:classification-kinship:attribute}{अभिलक्षण} साझा है? \omterm{prop:g6:classification-kinship:kinship}{नातेदारी} की भाषा में सबसे
        चौड़ी साझेदारी का क्या अर्थ है?
  \item ``चार \omterm{def:g1:my-body:parts}{उपांग}'' पर कौन-कौन से नमूनों के निशान लगते हैं? कबूतर और
        डॉल्फ़िन के निशान समझाओ --- \omterm{def:g6:classification-kinship:tree}{वृक्ष} के तर्क से डैने और मीन-पंख
        हैं क्या?
  \item ``पानी में रहता है'' तालिका में कहीं क्यों नहीं आता, जबकि
        ट्राउट और डॉल्फ़िन, दोनों वहीं रहते हैं? अध्याय का नियम बताओ।
\end{enumerate}

\medskip\noindent\textbf{भाग II --- \omterm{def:g6:classification-kinship:tree}{वृक्ष}।}
\begin{enumerate}[resume]
  \item \omterm{def:g6:classification-kinship:tree}{वृक्ष} खींचो: दोराहों का क्रम रखो --- \omterm{def:g4:classifying-animals:vertebrate}{रीढ़}, चार \omterm{def:g1:my-body:parts}{उपांग}, और फिर
        आवरण (\omterm{def:g1:animals-around-us:covering}{शल्क}, \omterm{def:g1:animals-around-us:covering}{पंख}, \omterm{def:g1:animals-around-us:covering}{रोम}) अपनी-अपनी शाखाओं पर --- और छहों को उन पर
        रखो।
  \item अपने \omterm{def:g6:classification-kinship:tree}{वृक्ष} से पढ़ो: दराज़ में चूहे का सबसे क़रीबी चचेरा कौन है,
        और ट्राउट का?
  \item मेंढक ``चार \omterm{def:g1:my-body:parts}{उपांग}'' से आगे खड़ा है पर आवरण वाले हर दोराहे से
        पहले। उसकी खुली, नम \omterm{def:g1:the-five-senses:sense}{त्वचा}
        (\cref{rem:g4:classifying-animals:confusion}) उसे किस चीज़ का
        निशान बनाती है, और उसकी शाखा कहाँ अलग होती है?
  \item सातवें नमूने की चिट गिर जाती है: सूखी शल्कदार \omterm{def:g1:the-five-senses:sense}{त्वचा}, चार \omterm{def:g1:my-body:parts}{उपांग},
        \omterm{def:g4:classifying-animals:vertebrate}{रीढ़}। उसे \omterm{def:g6:classification-kinship:tree}{वृक्ष} पर रखो और चिट देखे बिना उसका वर्ग बताओ।
\end{enumerate}

\medskip\noindent\textbf{भाग III --- अर्थ।}
\begin{enumerate}[resume]
  \item डॉल्फ़िन और ट्राउट धारा-रेखित बनावट साझा करते हैं; और डॉल्फ़िन
        तथा चूहा रोम-और-दूध वाले पूर्वज। \omterm{def:g6:classification-kinship:tree}{वृक्ष} पर कौन-सी जोड़ी ज़्यादा
        पास बैठती है, और \omterm{def:g6:classification-kinship:tree}{वृक्ष} को धारा-रेखित बनावट को अनदेखा क्यों
        करना ही चाहिए?
  \item ``चार \omterm{def:g1:my-body:parts}{उपांग}'' वाला दोराहा इतिहास के बारे में क्या दावा करता
        है: बहुत पहले की कौन-सी एक ही चीज़ से मेंढक, छिपकली, कबूतर,
        चूहा और डॉल्फ़िन, सब उतरे हैं?
  \item कक्षा का \omterm{def:g6:classification-kinship:tree}{वृक्ष} हर शाखा को आख़िरकार तने पर मिला देता है। एक
        वाक्य में बताओ कि वह चित्र इन छहों नमूनों के बारे में ---
        और उसे खींचने वाले कक्षा के सदस्यों के बारे में --- चुपचाप
        क्या दावा कर रहा है।
  \item एक शक्की कहता है: ``तुम्हारा \omterm{def:g6:classification-kinship:tree}{वृक्ष} तो बस फ़ाइल करने की एक
        व्यवस्था है।'' दो हिस्सों में उत्तर दो: \omterm{def:g6:classification-kinship:attribute}{अभिलक्षणों} का नत्थी
        होना किस चीज़ के बारे में एक \emph{तथ्य} है, और \omterm{def:g6:classification-kinship:tree}{वृक्ष} उस तथ्य
        की \emph{व्याख्या} क्या करता है जिसे महज़ फ़ाइल करना बिना
        समझाए छोड़ देता?
\end{enumerate}
\end{problem}
